密文与 $\Lambda$ 的整数提升使用相同的环和尺度。为得到所述模 $t$ 的解码明文，假设 $t\mid q$、$\Delta=q/t$ 且 $\|E\|_\infty<\Delta/2$。

假设有一个 GLWE 密文 \textsf{ct} 和一个新的明文多项式 $\Lambda$，如下所示：


$\textsf{ct} = \textsf{GLWE}_{S, \sigma}(\Delta M + E) = ( A_0, A_1, \ldots, A_{k-1}, B) \in \mathcal{R}_{\langle n,q \rangle}^{k + 1}$

$\Lambda$：新的明文多项式

$\Delta \Lambda$：按 $\Delta$ 放大后的新明文多项式

$ $

\noindent 定义如下密文与明文加法运算：

$\textsf{ct} + \Delta\Lambda = ( A_0, \text{ } A_1, \ldots, A_{k-1}, \text{ } B + \Delta \Lambda)$

$ $

\noindent 则有以下结论：

\begin{tcolorbox}[title={\textbf{\tboxlabel{\ref*{sec:glwe-add-plain}} GLWE 密文与明文的同态加法}}]
$\textsf{GLWE}_{S, \sigma}(\Delta M + E) + \Delta\Lambda $

$=  (\{A_i\}_{i=0}^{k-1}, \text{ } B) + \Delta\Lambda$

$=  (\{A_i\}_{i=0}^{k-1}, \text{ } B + \Delta\Lambda)$

$= \textsf{GLWE}_{S, \sigma}(\Delta (M + \Lambda) + E)$
\end{tcolorbox}


这意味着，将按 $\Delta$ 放大后的明文多项式 $\Lambda$ 加到加密 $M$ 的 GLWE 密文上，再解密所得密文，便得到 $M + \Lambda$。




\begin{myproof}

\begin{enumerate}
\item 由于 $B = \sum\limits_{i=0}^{k-1}{(A_i \cdot S_i)} + \Delta \cdot M + E$，

$B + \Delta \cdot \Lambda = \sum\limits_{i=0}^{k-1}{(A_i \cdot S_i)} + \Delta \cdot M + E + \Delta \cdot \Lambda = \sum\limits_{i=0}^{k-1}{(A_i \cdot S_i)} + \Delta \cdot (M + \Lambda) + E$  \\

这意味着 $(A_0, A_1, ... \text{ } A_{k-1}
, B + \Delta \Lambda)$ 构成密文 $\textsf{GLWE}_{S, \sigma}(\Delta (\Lambda + M) + E)$。
\item 因此， \\
$\textsf{GLWE}_{S, \sigma}(\Delta M + E) + \Delta\Lambda$ \\
$ = (A_0, \text { } A_1, ... \text{ } A_{k-1}, \text { } B + \Delta \Lambda)$ \\
$ = (\{A_i\}_{i=0}^{k-1}, \text { } B + \Delta \Lambda)$ \\
$= \textsf{GLWE}_{S, \sigma}(\Delta (M + \Lambda) + E)$

%\begin{flushright}
%\qedsymbol{} 
%\end{flushright}
\end{enumerate}

\end{myproof}


\para{噪声增长：}相对于未约减消息 $M+\Lambda$，相位误差仍为 $E$。当 $t\mid q$ 时，将消息模 $t$ 约减不会改变有效误差。若 $t\nmid q$，消息回绕可能引入 \autoref{subsubsec:lwe-noise-bound} 描述的溢出误差；此时“噪声不变”只能理解为相对于未约减消息而言。